% Generated draft: requires mathematical and source review.
\documentclass[11pt]{article}
\usepackage[margin=25mm]{geometry}
\usepackage{fontspec}
\setmainfont{Latin Modern Roman}
\usepackage{amsmath,amssymb,mathtools,unicode-math}
\setmathfont{Latin Modern Math}
\usepackage{enumitem,needspace,xurl,hyperref}
\setlength{\parindent}{0pt}
\setlength{\parskip}{5pt}
\setlength{\emergencystretch}{3em}
\newcommand{\N}{\mathbb{N}}
\newcommand{\Z}{\mathbb{Z}}
\newcommand{\Q}{\mathbb{Q}}
\newcommand{\R}{\mathbb{R}}
\newcommand{\C}{\mathbb{C}}
\newcommand{\F}{\mathbb{F}}
\newcommand{\A}{\mathbb{A}}
\newcommand{\PP}{\mathbb{P}}
\newcommand{\E}{\mathbb{E}}
\newcommand{\eps}{\varepsilon}
\newcommand{\dd}{\,\mathrm{d}}
\begin{document}
\section*{UnsolvedMath ID: 20002870}
\textbf{Source problem code:} catalogue entry
\subsection{Definitions and mathematical statement}
\textbf{Question (B. Stensones).} Let \(p(z,\bar z,w,\bar w)\) be a real-valued plurisubharmonic homogeneous polynomial of degree \(2k\) on \(\C^2\). If an algebraic curve \(\Sigma\) avoids zero and \(p|_\Sigma=0\), must \(\Sigma\) be one level set of a homogeneous holomorphic polynomial? The source does not specify irreducibility, ordinary versus bihomogeneity, positivity beyond plurisubharmonicity, or exclusion of pluriharmonic terms. The aid uses ordinary real homogeneity \(p(tz,tw)=t^{2k}p(z,w)\), \(t>0\), and affine algebraic curves. A nonempty level avoiding the origin must use nonzero level value and positive polynomial degree.

\textbf{Reported reducible obstruction.} Under those literal conventions take
\[
 p(z,w)=\operatorname{Im}(zw),\qquad
 \Sigma=\{zw=1\}\cup\{zw=2\}.
\]
The polynomial is nonzero, homogeneous quadratic and pluriharmonic, hence plurisubharmonic. The pure one-dimensional algebraic curve avoids the origin and lies in its zero set. The aid reports that it is not a single level \(\{H=c\}\) for any positive-degree homogeneous holomorphic \(H\). Constancy of \(H\) at a nonzero value on \(zw=1\) forces \(H\) to be a scalar power of \(zw\), whose nonzero values at \(zw=1\) and \(zw=2\) differ. This counterexample relies on allowing reducibility and pluriharmonic polynomials; it does not refute a strengthened version excluding them.

\textbf{Reported irreducible quadratic theorem.} If \(p\) is any nonzero plurisubharmonic homogeneous quadratic and \(\Sigma\) is an irreducible algebraic curve avoiding zero with \(p|_\Sigma=0\), the aid reports
\[
 \Sigma=\{H=c\},\qquad c\ne0,
\]
for a homogeneous holomorphic polynomial \(H\) of degree one or two. Nonzero \(p\) and irreducibility are hypotheses of this special theorem. No analogous all-degree conclusion is supplied.

\textbf{Reported higher-degree reduction.} At every smooth point of \(\Sigma\), its complex tangent is Levi-null for \(p\), since the restriction is constant and the Levi form is nonnegative. If \(\Sigma=\{F=0\}\) with irreducible \(F\in\C[z,w]\), then \(F(0)\ne0\). For the Euler field \(E=z\partial_z+w\partial_w\), the aid gives the criterion
\[
 \Sigma\text{ is a homogeneous holomorphic level}
 \quad\Longleftrightarrow\quad
 EF\text{ is constant modulo }(F).
\]
Equivalently \(EF-c\in(F)\) for some scalar. For \(d=\deg F\), this condition forces \(F=F_d+F(0)\) with homogeneous \(F_d\), hence the nonzero level \(F_d=-F(0)\). The aid does not prove that tangential Levi-nullity forces this Euler criterion for every higher-degree irreducible curve. Thus the missing general step, rather than the reducible counterexample, is the target of a clarified irreducible version. All supplied claims require independent verification.
\subsection{Short English statement}
The literal wording admits a reducible quadratic counterexample using a pluriharmonic polynomial. The aid reports a positive theorem for nonzero quadratic polynomials and irreducible curves. In higher degree it reduces the desired conclusion to an Euler-field condition, without proving that condition from plurisubharmonicity.
\subsection{Sources}
\begin{itemize}
\item Local UnsolvedMath registry, permanent ID \(20002870\), catalogue entry, \texttt{aim-several-complex-variables-notes.json}, \texttt{notes[27]}, workshop \emph{Emerging applications of complexity for CR mappings}, ``Plurisubharmonic polynomials'', Problem 1, B. Stensones: \url{https://aimath.org/WWN/crmappings/crmappings.pdf}.
\item Supplied research summary on the reducible obstruction, irreducible quadratic theorem, tangential Levi-nullity and Euler criterion. No independent authoritative theorem references are supplied.
\item Attempt 1: partial\_result, status at run ``unknown'', rigor ``proved'', low novelty confidence; registry status partially\_solved. The machine-generated aid is not peer reviewed and requires independent verification before citation; no PDF visit is claimed.
\end{itemize}
\end{document}
